\section{Robustness}\label{sec:res-robust}
\Cref{sec:res-checks} introduces these two checks, and \Cref{tab:checks} sums them up.

% src: S-PAPER-1, S-PAPER-2, S-PAPER-3, S-PAPER-6, S-PAPER-10
The two robustness checks in this appendix speak to two limits.
In R1, only \RobROneSplitN{} kept markets have a split payout and \RobROneMissingN{} is missing from the file.
Our label agrees with the resolution file on \RobROneAgreement{} of clean markets.
The Q1 rule result is the same without the unclear markets.
In R2, the price moves by at least \DeltaA{} in \RobRTwoShareDeltaAtLeastTen{} of markets.
This is the move between \pOne{} and the last trade.
So \pOne{} is not the price at the close.
Also \RobRTwoShareNoTrade{} of markets have no trade in the last day.

\subsection{R1: split payouts}\label{sec:app-r1}
% src: S-PAPER-9
Check R1 asks whether the Q1 result holds when markets without a clean payout are removed.
The resolution file gives each market a payout vector.
A market is clean when one outcome pays one and the other pays zero.
Of the kept markets, \RobROneCleanN{} are clean (\RobROneCleanShare).
\RobROneSplitN{} have a split payout (\RobROneSplitShare).
\RobROneOtherN{} are of another kind.
\RobROneMissingN{} is missing from the file.
Our label marks outcome 1 as the winner on \RobROneSplitYesOne{} of the split markets and outcome 0 on \RobROneSplitYesZero.
Our winner label agrees with the resolution file on \RobROneAgreement{} of the \RobROneNClean{} clean markets.
\RobROneNMismatch{} clean markets disagree.
\Cref{fig:r1} shows the payout classes and the calibration points with and without the unclear markets.
\Cref{tab:r1} puts the sensitivity row next to the original row.
We remove split, other, and mismatching markets, which is \RobROneNExcluded{} markets.
The Q1 rule result is the same in both rows: \RobROneRuleOriginal{}
The chi-square statistic moves from \QOneChiSq{} to \RobROneSensChiSq.

\begin{figure}[H]
\centering
\figorbox{figures/fig_r1_split_payout}{\linewidth}
\caption{Kept markets by payout class (log scale), and the Q1 calibration points with and without the unclear markets. Why: points that overlap mean the unclear markets do not move the result.}% Source: run \RobustRun.
\label{fig:r1}
\end{figure}

\begin{table}[H]
\centering
\caption{Q1 original row and R1 sensitivity row. Why: close rows mean the unclear markets do not drive the Q1 result.}% Source: runs \RunId{} and \RobustRun.
\label{tab:r1}
\small
\begin{tabular}{lrrrrr}
\toprule
Row & Markets & $\chi^2$ & $p$ & Mean sq. residual & Bins off (BH) \\
\midrule
Q1 original & \QOneN & \QOneChiSq & \QOneP & \QOneBrier & \QOneBinsBh \\
\midrule
Sensitivity & \RobROneSensN & \RobROneSensChiSq & \RobROneSensP & \RobROneSensBrier & \RobROneSensBinsBh \\
\bottomrule
\end{tabular}
\end{table}

\subsection{R2: price change before the close}\label{sec:app-r2}
% src: S-PAPER-9
Check R2 asks how far the price moves from $T-\DaysShort$ day to the last trade before $T$.
Call the mean price at the last timestamp before $T$ \pT.
The median absolute change from \pOne{} to \pT{} is \RobRTwoMedianAbsDelta.
The change is at least \DeltaA{} in \RobRTwoShareDeltaAtLeastTen{} of markets.
It is at least \DeltaB{} in \RobRTwoShareDeltaAtLeastTwentyFive{} and at least \DeltaC{} in \RobRTwoShareDeltaAtLeastFifty.
The price crosses \FavoriteLine{} in \RobRTwoShareCrossing{} of markets.
This compares the two end points and not the path.
\RobRTwoShareNoTrade{} of markets have no trade in the last day, and for them \pT{} equals \pOne.
\Cref{fig:r2} shows the distribution of the absolute change.

\begin{figure}[H]
\centering
\figorbox{figures/fig_r2_price_path}{\linewidth}
\caption{ECDF of the absolute price change from \pOne{} to \pT. Why: a curve near one at small values means \pOne{} is close to the last price.}% Source: run \RobustRun.
\label{fig:r2}
\end{figure}

The price \pT{} is the price just before the close.
So its mean squared residual is not a one-day-ahead forecast.
The mean squared residual is \RobRTwoAllBrierOneDay{} at \pOne{} and \RobRTwoAllBrierAtClose{} at \pT.
The paired difference (\pT{} minus \pOne) is \RobRTwoAllMeanDiff{}
(\CiLevel{} CI [\RobRTwoAllCiLow, \RobRTwoAllCiHigh], $p\,\RobRTwoAllP$).
A negative difference means \pT{} is closer to the outcome than \pOne.
On the \RobRTwoTradedN{} markets with a trade in the last day, the difference is \RobRTwoTradedMeanDiff.
The \CiLevel{} CI is [\RobRTwoTradedCiLow, \RobRTwoTradedCiHigh].

% src: owner, 2026-11-18 (D-053); numbers from run \RTwoCloseRun
We also repeat the Q1 calibration with \pT{} in place of \pOne.
This check is descriptive and was not pre-registered.
It uses the same markets, bins, and tests as Q1.
With \pOne, the same code gives the Q1 values ($\chi^2=\RTwoCloseOneDayChiSq$).
With \pT, the chi-square statistic is \RTwoCloseCloseChiSq{} with \QOneDf{} degrees of freedom ($p=\RTwoCloseCloseP$).
So the chi-square test still rejects.
In \RTwoCloseCloseBinsNeg{} of \GapBinsTotal{} bins the observed rate is below the mean price.
The BH correction flags \RTwoCloseCloseBinsBh{} bins, and Bonferroni flags \RTwoCloseCloseBinsBonferroni.
The net gap shrinks from \RTwoCloseOneDayNetGap{} to \RTwoCloseCloseNetGap.
So the timing of the price explains part of the gap, but not all of it.
Near the close, many markets move toward 0 or 1.
For example, bin 0 grows from \RTwoCloseOneDayBinZeroN{} to \RTwoCloseCloseBinZeroN{} markets.
So a bin does not hold the same markets at the two prices.
\Cref{fig:r2-close} shows the per-bin gaps at both prices.

\begin{figure}[H]
\centering
\figorbox{figures/fig_r2_close_gap}{5.5in}
\caption{Per-bin gap (observed rate minus mean price) at \pOne{} and at \pT, with \CiLevel{} intervals. Why: if the last day's news closed the gap, the blue points would sit on the zero line.}% Source: run \RTwoCloseRun.
\label{fig:r2-close}
\end{figure}
